\chapter{Why the generated algebra is rational}
\label{ch:regularity}

The construction so far gives an infinite-dimensional unitary algebra
inside a lattice theory. Rationality asks for a much stronger form of
control: a finite and semisimple representation theory. A unitary
vacuum representation does not automatically provide this. The
positive conformal-vector identity from the previous chapter is the
extra ingredient. We will use it first to bound a quotient of the
vacuum algebra, then to constrain arbitrary modules, and finally to
split extensions of the vacuum. A categorical theorem turns the last
property into semisimplicity.

The proof is worth following even for readers primarily interested in
physics. It explains why the later finite list of characters is a
conclusion of the construction rather than an initial assumption.

\section{The finiteness properties that must be distinguished}
\index{regularity!vertex algebras}\index{C2-cofiniteness@$C_2$-cofiniteness}

For a VOA $V$ define
\[
 C_2(V)=\Span\{a_{(-2)}b:a,b\in V\}.
\]
The notation $C_2$ here has no relation to central charge two. The VOA
is $C_2$-cofinite when $V/C_2(V)$ is finite dimensional.
This quotient removes many derivatives and normal composites, so it
measures a different kind of finiteness from $\dim V_n<\infty$.
Finite-dimensional homogeneous spaces alone do not imply
$C_2$-cofiniteness.

\begin{example}[The Heisenberg warning]
A rank-one Heisenberg VOA has finite-dimensional weight spaces, but
the classes of $J_{(-1)}^n\one$ survive in its $C_2$ quotient. The
quotient is a polynomial algebra in the current class and is infinite
dimensional. Correspondingly, there are continuously many Fock modules
labeled by momentum. A simple positive vacuum representation is fully
compatible with this nonrational behavior.
\end{example}

\begin{example}[A five-dimensional quotient of the rank-one lattice theory]
\label{ex:regularity-A1-C2}
\index{C2-cofiniteness@$C_2$-cofiniteness!rank-one lattice example}
Let $V=V_{A_1}=L_1(\mathfrak{sl}_2)$, with the weight-one
currents $E,F,H$ normalized as in
Example~\ref{ex:c-root-triple}.
The VOA identities make
\[
 R=V/C_2(V),\qquad [a][b]=[a_{(-1)}b],
\]
a commutative associative graded algebra. Derivatives vanish in
$R$, so strong generation by $E,F,H$ implies that their classes
$e,f,h$ generate it as an ordinary commutative algebra.

The lattice exponential formula supplies the decisive relation.
For a root $\alpha$ of squared norm two, $Y(e^\alpha,z)e^\alpha$
starts at $z^2$. Its constant coefficient is therefore zero:
\[
 E_{(-1)}E=0,\qquad e^2=0.
\]
This is a relation in the actual simple lattice VOA. The universal
affine vacuum algebra at level one has not yet imposed this
integrability relation.

Zero modes preserve $C_2(V)$ and act as derivations on $R$.
The derivation $\partial_F$ induced by $F_{(0)}$ obeys
\[
 \partial_F(e)=-h,\qquad
 \partial_F(h)=2f,\qquad \partial_F(f)=0.
\]
Apply it successively to $e^2=0$, $eh=0$, $h^2-2ef=0$,
and $hf=0$. The resulting equations are
\[
 -2eh=0,\qquad -h^2+2ef=0,\qquad
 6hf=0,\qquad 2f^2=0.
\]
Consequently every class is a linear combination of
\[
 1,\quad e,\quad h,\quad f,\quad ef.
\]
Every cubic monomial vanishes; for instance
$h^3=2ehf=0$. This proves an upper bound of five, but it does
not yet prove that the last class survives.

Weights zero and one already give four independent classes,
since $C_2(V)$ has no component in either weight. At weight two,
the weight formula for $a_{(-2)}b$ gives exactly
\[
 C_2(V)\cap V_2=D V_1,\qquad D=L_{-1}.
\]
The conformal vector satisfies $L_1\omega=0$, so
$\langle Da,\omega\rangle=\langle a,L_1\omega\rangle=0$
for $a\in V_1$. Its squared norm is $c/2=1/2$, hence
$\omega\notin C_2(V)$. Finally the level-one Sugawara formula
gives
\[
 [\omega]=\frac16\left(\frac{h^2}{2}+2ef\right)
          =\frac{ef}{2}=\frac{h^2}{4}.
\]
Thus $ef\ne0$, and the displayed spanning set is a basis:
\[
 \dim R=5,\qquad
 \sum_{n\geq0}(\dim R_n)t^n=1+3t+t^2.
\]

Adding the lattice exponential fields to the Heisenberg algebra
has made the surviving current classes nilpotent and the quotient
finite. This establishes $C_2$-cofiniteness; by itself it is not
a proof of rationality. Nor is five a count of simple modules:
this VOA has the two affine simples labelled $\ell=0,1$.
\end{example}

An ordinary module has diagonalizable $L_0$, finite-dimensional
eigenspaces, and lower truncation in each integer translate of an
eigenvalue. The simple ordinary modules used below have spectrum in
one $h+\Z_{\geq0}$. A grading-restricted generalized module
allows finite Jordan blocks for $L_0$. Such modules have to be
considered before rationality has been proved: excluding them at the
definition stage would assume part of the conclusion.

For the present discussion, rationality means semisimplicity of the
appropriate grading-restricted module theory; under the standard
$C_2$-cofinite hypotheses this agrees with the usual rationality
condition. A VOA is strongly rational when it is simple, of CFT type,
self-contragredient, rational, and $C_2$-cofinite. Regularity of the
component algebras below means that every weak module decomposes as
an algebraic direct sum of ordinary simple modules. This stronger
component property is known for the positive-integral affine,
Potts, and parafermion algebras being used.

\section{Compatible gradings and partial energies}

Let $A\subset V$ be a vertex subalgebra with its own conformal vector
$e$. We do not require $e=\omega_V$. The inclusion is
\emph{grading compatible} when a homogeneous state of $A$ has the
same intrinsic weight in $A$ as its ambient weight in $V$.
The two energy operators need not agree on all of $V$.

This distinction is familiar for a tensor product: $A\otimes\one$
has its own stress and its own energy, while the total energy also
counts the other factor. The Haagerup construction is subtler because
the many regular subalgebras in its positive identity need not commute.
Nevertheless each inclusion is grading compatible, so a mode of an
$A$-state changes ambient and intrinsic degrees by the same amount.

We use the following precise package of hypotheses.

\begin{definition}[A positive regular stress family]
A simple unitary VOA $V$ of CFT type has a positive regular stress
family if there are finitely many PCT-stable, grading-compatible, regular
unitary subalgebras, carrying the inherited positive form,
$A_i\subset V$ with conformal vectors $e_i$ such that
\begin{equation}\label{eq:c-regular-positive}
 \omega_V=\sum_i t_i e_i,\qquad t_i>0,
\end{equation}
and each $A_i$ has only nonnegative simple lowest weights, with the
vacuum its unique simple module of lowest weight zero. We require
each $e_i$ to be an ambient weight-two quasiprimary fixed by
PCT:
\[
 L^V_0e_i=2e_i,\qquad L^V_1e_i=0,\qquad\theta(e_i)=e_i.
\]
The invariant-form identity then gives
$(L^{e_i}_n)^\dagger=L^{e_i}_{-n}$ on $V$.
\end{definition}

The affine, Potts, and PF$_{13}$ embeddings established in the preceding
chapter give exactly such a family for $U$. Their PCT-real stresses
are ambient quasiprimary vectors of weight two; the invariant-form
adjoint identity therefore gives the displayed mode adjoints.
In particular the partial energies are self-adjoint on each ambient
grade. No properness assumption
is needed here. The regularity argument would still work if the
generated algebra had turned out to be the ambient $B$.

\section{\texorpdfstring{Bounding the $C_2$ quotient by energy}{Bounding the C2 quotient by energy}}

Regularity implies $C_2$-cofiniteness by Li's theorem; this is an
available implication, not a further assumption about the embedded
components. We also use the module-cofiniteness theorem for a regular
$C_2$-cofinite algebra $A$: for each ordinary simple $A$-module $M$,
the quotient
\[
 M/C_{2,A}(M),\qquad
 C_{2,A}(M)=\Span\{a_{(-2)}m:a\in A,\ m\in M\},
\]
is finite dimensional. The theorem includes modules, not only the
vacuum algebra. Because $A$ has finitely many simple modules, one can
choose a uniform upper bound on the intrinsic energies that survive
in these quotients.

\begin{proposition}[Positive stresses imply $C_2$-cofiniteness]
\label{prop:c-c2}
A VOA with a positive regular stress family is $C_2$-cofinite.
\end{proposition}
\begin{proof}
For each $A_i$, choose a number $D_i$ bounding the intrinsic energy
of every quotient $M/C_{2,A_i}(M)$ over its finitely many simple
modules. Regularity decomposes the weak $A_i$-module $V$ into such
modules. The component $C_2$ construction commutes with an algebraic
direct sum:
$C_{2,A_i}(\bigoplus_jM_j)=\bigoplus_jC_{2,A_i}(M_j)$.
An isomorphism of component modules preserves their intrinsic
conformal operator, so repeated copies do not introduce arbitrary
intrinsic energy shifts. Thus the same $D_i$ works even when a
component simple occurs infinitely often in $V$.
Therefore every $L^{e_i}_0$ eigenvector of energy greater
than $D_i$ belongs to $C_{2,A_i}(V)$, which is contained in $C_2(V)$.

The operator $L^{e_i}_0$ preserves each finite-dimensional ambient
space $V_n$ and is self-adjoint for its positive form. Take
$v\in V_n\cap C_2(V)^\perp$. Its components along all intrinsic
energies greater than $D_i$ vanish, so
\[
 \langle L^{e_i}_0v,v\rangle\leq D_i\|v\|^2.
\]
Taking the zero mode of \eqref{eq:c-regular-positive} yields
\[
 n\|v\|^2=\sum_i t_i\langle L^{e_i}_0v,v\rangle
                 \leq\left(\sum_i t_iD_i\right)\|v\|^2.
\]
Thus no nonzero $v$ exists once $n>\sum_i t_iD_i$. Since $C_2(V)$
is ambient graded, these orthogonal complements represent the graded
pieces of $V/C_2(V)$. Only finitely many finite-dimensional grades
survive.
\end{proof}

Only one partial energy is diagonalized at a time. The proof would be
invalid if it assumed a simultaneous eigenbasis for all $L^{e_i}_0$.
It instead applies separate expectation-value bounds to the same vector
and adds them using positive coefficients. This is the precise point
of the positive-stress construction.

The proof is existential: it gives a finite cutoff from the component
quotients but does not compute a sharp numerical cutoff for $U/C_2(U)$.
That is enough for the structural conclusion. A tight cutoff would be
useful for an explicit presentation of its Zhu algebra, the associative
algebra that encodes degree-zero mode actions on lowest spaces. Such a
presentation is not required to establish rationality by the route below.

\section{Lowest spaces are complete component tops}

Before applying positivity to arbitrary $V$-modules we must avoid an
easy mistake. A partial conformal energy can be positive on the vacuum
Hilbert space without a positive form having yet been supplied on every
module. The argument below uses regular component representation theory
and traces; it does not assume all $V$-modules are already unitary.

For a weak $A$-module $M$, its highest space is
\[
 \Omega_A(M)=\{v:a_{(n)}v=0\text{ for homogeneous }a\in A,
                                  \ n\geq\wt_A(a)\}.
\]
These are precisely the modes that lower intrinsic degree. The
degree-zero modes $o(a)=a_{(\wt(a)-1)}$ act on this space through
the Zhu algebra. Regularity makes that finite-dimensional top algebra
semisimple on the highest spaces under consideration.

\begin{lemma}[Compatible top spaces]\label{lem:c-compatible-tops}
Let $A\subset V$ be regular and grading compatible. A finite space
$E$ in a weak $V$-module, stable under all $A$ degree-zero modes and
killed by all negative-degree $A$ modes, is a direct sum of simple
$A$ top representations. In particular, $L^A_0$ is diagonalizable on
$E$ with eigenvalues among the simple lowest weights of $A$.
\end{lemma}
\begin{proof}
Compatibility places $E$ in $\Omega_A(M)$. Decompose $M$ as a weak
$A$-module using regularity. The highest space of an ordinary simple
module is its top: a vector of higher energy killed by every lowering
mode would, after mode reordering, generate a nonzero submodule
at its own energy and above. That submodule could not contain the
original lowest space, contradicting simplicity.
Thus $\Omega_A(M)$ is a sum of those tops. The
finite invariant subspace $E$ is consequently a sum of simple
representations of the semisimple top algebra.
\end{proof}

The condition concerns lowering modes of \emph{all} $A$-states.
Checking only a few generator modes requires an additional argument
that they imply the full highest-space condition.

\begin{proposition}[Nonnegative lowest weights]
Every simple ordinary $V$-module for a positive regular stress family
has real nonnegative lowest weight. Its only simple module of lowest
weight zero is the vacuum module $V$.
\end{proposition}
\begin{proof}
Let $E=M_h$ be the finite lowest space. Degree-zero modes preserve it
and lowering modes kill it. By Lemma~\ref{lem:c-compatible-tops}, each
$L^{e_i}_0$ is diagonalizable on $E$ with nonnegative real spectrum.
Taking traces gives
\begin{equation}\label{eq:c-top-positive-trace}
 h\dim E=\sum_i t_i\Tr_E L^{e_i}_0\geq0.
\end{equation}
This proves the assertion about $h$ even if the partial energies do
not commute.

If $h=0$, all nonnegative traces vanish. Each partial energy is
diagonalizable, so it vanishes on $E$. The unique-zero-top hypothesis
puts $E$ in sums of vacuum tops for every $A_i$. Therefore
$L^{e_i}_{-1}E=0$, and the conformal-vector identity gives
$L^V_{-1}E=0$. The vacuum-like lemma below supplies a nonzero module
map $V\to M$. Simplicity makes it an isomorphism.
\end{proof}

\section{A translation-invariant vector is a vacuum vector}

\begin{lemma}[Vacuum-like vectors]\label{lem:c-vacuumlike}
\index{vacuum-like vector}
If $v$ in a weak $V$-module satisfies $L_{-1}v=0$, then
$a_{(n)}v=0$ for every $a\in V$ and $n\geq0$. The map
\[
 s:V\longrightarrow M,\qquad s(a)=a_{(-1)}v
\]
is a module homomorphism with $s(\one)=v$.
\end{lemma}
\begin{proof}
Lower truncation says that $a_{(n)}v=0$ for sufficiently large $n$.
Translation covariance gives
\[
 L_{-1}a_{(n+1)}v=-(n+1)a_{(n)}v.
\]
Downward induction proves vanishing for every $n\geq0$. Taking the
appropriate constant coefficient in module associativity, or using
the module Jacobi identity, now gives
$b_{(m)}(a_{(-1)}v)=(b_{(m)}a)_{(-1)}v$. This is precisely the
homomorphism assertion.
\end{proof}

This lemma explains why the translation mode, rather than only energy
zero, appears in the projectivity proof. A generalized weight-zero
vector could initially belong to a nontrivial Jordan extension. The
positive component structure forces it to satisfy the full vacuum
condition.

\section{Splitting every extension onto the vacuum}

In an abelian category, an object $P$ is projective when every
epimorphism onto $P$ splits, equivalently when $\Hom(P,-)$ is exact.
Think of a possible nonsplit extension
\[
 0\longrightarrow X\longrightarrow E\longrightarrow V
                                      \longrightarrow0.
\]
Projectivity says that a vacuum copy can always be lifted back into
$E$ as an actual submodule.

\begin{proposition}[The vacuum is projective]\label{prop:c-projective}
\index{projective vacuum}
For a positive regular stress family, $V$ is projective in the
category of grading-restricted generalized modules.
\end{proposition}
\begin{proof}
Let $\pi:E\twoheadrightarrow V$. The standard $C_2$-cofinite module
theory gives finite length and makes its simple composition factors
ordinary modules. The preceding lowest-weight result therefore implies
that every composition factor has nonnegative energy. Generalized
weight spaces are exact in these short exact sequences: module maps
commute with $L_0$, and every vector has a finite decomposition into
generalized eigenspaces. Induction along a composition series then
shows that $E$ has no negative generalized weights. Exactness also
gives a surjection $E_{[0]}\twoheadrightarrow V_0$.

Consider the finite generalized weight-zero space $E_{[0]}$. It is
killed by every negative-degree mode and is invariant under
degree-zero modes. By compatible-top semisimplicity, each partial
energy $L^{e_i}_0$ is diagonalizable there with nonnegative spectrum.
But $L^V_0$ is nilpotent on $E_{[0]}$, so
\[
 0=\Tr_{E_{[0]}}L^V_0
       =\sum_i t_i\Tr_{E_{[0]}}L^{e_i}_0.
\]
Every partial energy vanishes. As before, their vacuum-top translation
operators vanish, so $L^V_{-1}E_{[0]}=0$.

Since $\pi$ respects generalized weights, choose $v\in E_{[0]}$
with $\pi(v)=\one$. Lemma~\ref{lem:c-vacuumlike} gives
$s(a)=a_{(-1)}v$. Then $\pi s(a)=a_{(-1)}\one=a$, so $s$ is a
section of $\pi$.
\end{proof}

No Hermitian form on $E$ was used. The positivity in the proof lies
in the known nonnegative spectra of the regular component tops and
the positive coefficients of the stress identity. This avoids
assuming strong unitarity before it is proved by the later analytic
argument.

\section{From a projective unit to a semisimple category}

At this stage standard $C_2$-cofinite VOA tensor theory provides a
finite complex linear abelian braided category $\mathcal C$ of
generalized modules. Its tensor product is right exact. Contragredient
duality is a Grothendieck--Verdier duality with dualizing object $V$,
which is the tensor unit. These are the categorical inputs to the
next argument; we do not assume rigidity or semisimplicity already.

\index{Grothendieck--Verdier duality}
In this categorical section $D$ denotes contragredient duality, so $DX$
is the contragredient object; here $D$ is not the translation operator.
More precisely, this duality gives natural isomorphisms
\[
 \Hom(Y\otimes X,\one)\simeq\Hom(Y,DX).
\]
At $Y=DX$, the identity of $DX$ corresponds to the evaluation
$DX\otimes X\to\one$, so this map is nonzero when $X\ne0$.
The assertion represents maps \emph{into} the unit; it does not yet
supply a coevaluation from the unit satisfying the snake identities
or the exactness consequences of rigidity.

\begin{example}[Right exactness already bounds growth]
Let $A$ be a finite-dimensional algebra, and let $P=A$ be its
left regular module. A module $M$ of vector-space dimension $n$
is a quotient of $P^{\oplus n}$: choose a vector-space basis
$m_1,\ldots,m_n$ and send the $i$th copy's unit to $m_i$.
For any additive right-exact functor $G$ between finite module
categories, this gives a surjection
$G(P)^{\oplus n}\twoheadrightarrow G(M)$, hence
$\dim G(M)\leq n\dim G(P)$.

This estimate does not require $G$ to preserve injections.
For instance, with $A=\C[\varepsilon]/(\varepsilon^2)$,
$G(M)=M/\varepsilon M$ is right exact. The inclusion
$\varepsilon A\hookrightarrow A$ becomes the zero map
$\C\to\C$, so $G$ is not left exact. The dimension bound
still holds. This is the modest amount of exactness used in
the growth argument below.
\end{example}

We invoke the following rigidity theorem in its precise form: in a
Cauchy-complete linear braided category with finite-dimensional Hom
spaces and $\End(\one)=\C$, a nonzero indecomposable $X$ is rigid
if (i) $\one$ is a direct summand of $Y\otimes X$ for some $Y$, and
(ii) $\dim\End(X^{\otimes n})<n!$ for some $n$. These are called
non-negligibility and moderate growth. The theorem does not assume
tensor exactness.
The finite abelian category at hand is Cauchy complete, meaning that
every idempotent splits as a projection onto an actual direct summand.
It also has finite-dimensional Hom spaces and finite decompositions
into indecomposables, so the theorem can be applied one summand at a time.

\begin{proposition}[Projective unit implies semisimplicity]
\index{rigidity!from a projective unit}
\label{prop:c-projective-unit}
Let $\mathcal C$ be finite complex linear abelian and braided, with
simple unit, right-exact tensor product, and Grothendieck--Verdier
duality with dualizing object $\one$. If $\one$ is projective, then
$\mathcal C$ is rigid and semisimple.
\end{proposition}
\begin{proof}
Realize the underlying finite abelian category as finite-dimensional
modules over a finite-dimensional algebra. Let $P$ correspond to its
regular module and let $d_0(Z)$ denote the underlying vector-space
dimension. Every $Z$ is a quotient of $P^{\oplus d_0(Z)}$.
For fixed $X$, right exactness gives
\[
 d_0(Z\otimes X)\leq C_Xd_0(Z),\qquad
 C_X=\max(1,d_0(P\otimes X)).
\]
Iterating,
\[
 \dim\End(X^{\otimes n})
 \leq d_0(X^{\otimes n})^2
 \leq d_0(\one)^2C_X^{2n}<n!
\]
for sufficiently large $n$. The last inequality is the elementary
fact that factorial growth eventually dominates any fixed exponential.

For nonzero indecomposable $X$, duality gives a nonzero evaluation
$DX\otimes X\to\one$. Since the unit is simple, it is surjective.
Projectivity makes it split. The rigidity theorem therefore applies
to $X$. Finite direct sums and Krull--Schmidt decompositions make all
objects rigid.

Only now may we use that tensoring is exact. Rigidity identifies
\[
 \Hom(X,-)\cong\Hom(\one,-\otimes X^*).
\]
The right side is exact because the unit is projective and tensoring
by $X^*$ is exact. Thus every object is projective. Every short exact
sequence splits, which is semisimplicity.
\end{proof}

The order is important. Using tensor exactness to prove rigidity would
create a circular proof. The growth bound above uses only right
exactness, a property available before rigidity.

To connect this categorical conclusion with ordinary VOA rationality,
we use one more standard consequence of $C_2$-cofiniteness: every
finitely generated weak module belongs to the grading-restricted
generalized category. Each vector of any weak module belongs to its
cyclic submodule. Once that finite category is semisimple, these cyclic
submodules are sums of ordinary simple modules. The whole weak module
is therefore a sum of simple submodules, hence a direct sum of them
by the usual semisimple-module argument. In this setting the conclusion
actually gives regularity as well as rationality.

\begin{theorem}[Rationality of the Haagerup candidate]
Under the actual embedded-factor and positive-stress inputs of
Chapter~\ref{ch:factors}, the generated $U$ is strongly rational.
\end{theorem}
\begin{proof}
The lattice construction gives simplicity, CFT type, and a positive
invariant form; the latter gives self-contragredience. The positive
stress family gives $C_2$-cofiniteness by
Proposition~\ref{prop:c-c2} and projectivity of the vacuum by
Proposition~\ref{prop:c-projective}. The standard finite VOA category
and Proposition~\ref{prop:c-projective-unit} give semisimplicity,
hence rationality.
\end{proof}

This theorem does not identify the module category with the Haagerup
center. It does not yet construct positive forms on all simple modules
or on their intertwining-operator spaces. What it does provide is the
finite semisimple setting in which the branching and character arguments
of the next chapters become legitimate.

\section{Exercises}

\begin{exercise}\label{ex:c-regular-1}
Explain why $C_2(V)$ is homogeneous for ambient weight. Show that the
orthogonal complement in each $V_n$ represents the grade-$n$ quotient
even if $C_2(V)$ is not closed in the Hilbert completion.

\par\smallskip\noindent{\footnotesize\normalfont
\hyperref[sol:ex:c-regular-1]{Solution on page~\pageref*{sol:ex:c-regular-1}.}}\par
\end{exercise}
\begin{exercise}\label{ex:c-regular-2}
Prove that a diagonalizable operator with nonnegative real eigenvalues
and zero trace is zero. Give a nondiagonalizable counterexample if
diagonalizability is omitted.

\par\smallskip\noindent{\footnotesize\normalfont
\hyperref[sol:ex:c-regular-2]{Solution on page~\pageref*{sol:ex:c-regular-2}.}}\par
\end{exercise}
\begin{exercise}\label{ex:c-regular-3}
In the vacuum-like lemma, write the downward induction explicitly for
$a_{(2)}v$, $a_{(1)}v$, and $a_{(0)}v$. Why does the coefficient
$n+1$ never obstruct the induction through $n=0$?

\par\smallskip\noindent{\footnotesize\normalfont
\hyperref[sol:ex:c-regular-3]{Solution on page~\pageref*{sol:ex:c-regular-3}.}}\par
\end{exercise}
\begin{exercise}\label{ex:c-regular-4}
For the algebra $\C[\varepsilon]/(\varepsilon^2)$, exhibit a nonsplit
short exact sequence of modules with one-dimensional endpoints.
Compare the role of a projective simple module with the vacuum
projectivity argument above.

\par\smallskip\noindent{\footnotesize\normalfont
\hyperref[sol:ex:c-regular-4]{Solution on page~\pageref*{sol:ex:c-regular-4}.}}\par
\end{exercise}
\begin{exercise}\label{ex:c-regular-5}
Prove directly that $A^n<n!$ for all sufficiently large $n$, for any
fixed $A>0$. Use it to fill in the moderate-growth bound.

\par\smallskip\noindent{\footnotesize\normalfont
\hyperref[sol:ex:c-regular-5]{Solution on page~\pageref*{sol:ex:c-regular-5}.}}\par
\end{exercise}
\begin{exercise}\label{ex:c-regular-6}
List the point in the proof at which each of these hypotheses is used:
positive stress coefficients; component regularity; grading
compatibility; the unique component vacuum at weight zero; braiding;
right exactness. Which arguments would fail if a component had a
second simple module of weight zero?

\par\smallskip\noindent{\footnotesize\normalfont
\hyperref[sol:ex:c-regular-6]{Solution on page~\pageref*{sol:ex:c-regular-6}.}}\par
\end{exercise}
